\section{Worst-Case Convergence Analysis of the OMEGA Clause-Feedback Rule}

This section analyzes the stochastic clause-feedback dynamics used by the
current OMEGA Boolean-cell prototype.  Let $p\in(0,1)$ denote the exploration
probability.  At each non-satisfying state the algorithm selects a violated
clause; with probability $p$ it flips a uniformly selected variable of that
clause, and otherwise it chooses a variable minimizing the one-flip energy
change.

\subsection{Almost-Sure Convergence on Satisfiable Instances}

Fix any satisfiable 3-CNF formula $F$ and one satisfying assignment $x^\star$.
At a non-satisfying state $x$, every violated clause contains at least one
variable on which $x$ differs from $x^\star$.  During an exploration step, the
probability of selecting such a variable is at least $1/3$.  Hence the
conditional probability of decreasing the Hamming distance
$d_H(x,x^\star)$ by one is at least
\begin{equation}
q=\frac{p}{3}.
\end{equation}
Consequently, from any state, the probability of reaching $x^\star$ during the
next at most $n$ updates by a sequence of distance-decreasing exploration
moves is at least
\begin{equation}
q^n=\left(\frac{p}{3}\right)^n.
\end{equation}
It follows that the probability of never reaching a satisfying assignment is
zero.  A direct geometric-block estimate gives
\begin{equation}
\mathbb{E}[T]
\le
n\left(\frac{3}{p}\right)^n.
\end{equation}
Thus the current stochastic dynamics is globally convergent with probability
one on satisfiable instances, but this general bound is exponential.

\subsection{A Polynomial-Size Barrier Family}

The stronger desired statement---worst-case polynomial-time global
convergence---is false for the present update rule.  Consider the family
\begin{equation}
F_n
=
(x_n)
\land
\bigwedge_{i=1}^{n-1}
\left[
(\neg x_i\lor x_{i+1})^{w_i}
\land
(x_i\lor\neg x_{i+1})^{w_i}
\right],
\qquad
w_i=n-i.
\end{equation}
Here the superscript denotes repeated copies of the clause.  Unit and
two-literal clauses may be padded to three literals by repeating a literal,
without changing their truth value or the implementation's candidate-variable
set.  The number of clauses is
\begin{equation}
m
=
2\sum_{i=1}^{n-1}(n-i)+1
=
n(n-1)+1
=
\Theta(n^2).
\end{equation}
The formula has the unique satisfying assignment
\begin{equation}
x^\star=(1,1,\ldots,1).
\end{equation}

Define the suffix states
\begin{equation}
S_k
=
0^{\,n-k}1^k,
\qquad
0\le k\le n.
\end{equation}
For the exact violated-clause energy $E_F$,
\begin{equation}
E_F(S_0)=1,
\qquad
E_F(S_k)=k
\quad(1\le k\le n-1),
\qquad
E_F(S_n)=0.
\end{equation}

At $S_k$, for $1\le k\le n-2$, the only violated clauses lie on the
boundary between the zero prefix and the one suffix.  Flipping the boundary
one back to zero moves to $S_{k-1}$; flipping the boundary zero forward to one
moves to $S_{k+1}$.  The greedy branch therefore moves backward, because
\begin{equation}
E_F(S_{k-1}) < E_F(S_{k+1})
\end{equation}
(with the same conclusion at $k=1$, where the backward move has zero energy
change while the forward move increases the energy).

Let
\begin{equation}
a=\frac{p}{2},
\qquad
b=1-\frac{p}{2}.
\end{equation}
For $1\le k\le n-2$, the probability of advancing from $S_k$ to $S_{k+1}$
is $a$, while the probability of retreating to $S_{k-1}$ is $b$.  Since
$b>a$ for every $p<1$, the suffix length performs a biased random walk away
from the satisfying state.

Let
\begin{equation}
R=\frac{b}{a}
=
\frac{1-p/2}{p/2}>1.
\end{equation}
The gambler's-ruin probability that a walk beginning at $S_1$ reaches
$S_{n-1}$ before returning to $S_0$ is
\begin{equation}
h_n
=
\frac{R-1}{R^{n-1}-1}.
\end{equation}
Every visit to $S_0$ is followed by a forced move to $S_1$, so the expected
number of excursions required merely to reach $S_{n-1}$ is
\begin{equation}
\frac{1}{h_n}
=
\frac{R^{n-1}-1}{R-1}
=
\exp(\Omega(n)).
\end{equation}
This is already a lower bound on the expected time to reach the satisfying
state.

For the implementation value $p=0.45$,
\begin{equation}
R
=
\frac{0.775}{0.225}
=
3.444\ldots.
\end{equation}
Therefore the expected-excursion lower bound grows approximately as
$3.444^n$.  Because the formula itself has only $\Theta(n^2)$ clauses, this is
superpolynomial in the input size.

If the initial Boolean-cell state is uniformly random, the event $S_0=0^n$
has probability $2^{-n}$.  Hence even the unconditional expected runtime of
the uncapped single-run process obeys the lower bound
\begin{equation}
\mathbb{E}[T]
\ge
2^{-n}
\frac{R^{n-1}-1}{R-1}
=
\Omega\!\left(
\left(\frac{R}{2}\right)^n
\right).
\end{equation}
For $p=0.45$, $R/2=1.722\ldots>1$, so this lower bound is exponential in $n$.

\begin{theorem}
The present OMEGA clause-feedback rule with fixed exploration probability
$p=0.45$ does not have worst-case polynomial expected global-convergence time
over all formulas accepted by the implementation.
\end{theorem}

\begin{proof}
The family $F_n$ above has polynomial description size
$m=\Theta(n^2)$, but from $S_0$ the expected number of excursions required to
cross the energy barrier is
\[
(R^{n-1}-1)/(R-1)=\exp(\Omega(n)).
\]
Under uniform random initialization, multiplication by the probability
$2^{-n}$ of starting at $S_0$ still yields an exponential lower bound because
$R/2>1$ for $p=0.45$.  Hence the expected runtime is not bounded by any
polynomial in the input size.
\end{proof}

\subsection{Implication for the P=NP Program}

The exact SAT energy, sparse QUBO compilation, and OMEGA Boolean-cell
representation remain valid.  What fails is the specific local
clause-feedback dynamics as a candidate proof of polynomial global
optimization.  A successful replacement must introduce a genuinely nonlocal
progress mechanism whose worst-case number of updates can be bounded by a
polynomial; a local greedy/exploratory one-bit rule is insufficient.
